\solchapter{qft}

\sol{qft}{1} $\hbar c/mc^2$: electron $386\un{fm}$ ($3.9\times10^{-13}\un{m}$), the scale below which QED's pair creation matters and the electron cannot be localized as a single particle; proton $0.21\un{fm}$, a few times smaller than the proton's actual size ($0.84\un{fm}$, set by the strong force instead); $W$ boson $2.5\times10^{-3}\un{fm}$, the range of the weak force.

\sol{qft}{2} In natural units the Klein--Gordon equation is $(\partial_t^2 - \nabla^2 + m^2)\phi = 0$. Acting on the plane wave, $\partial_t\to-iE$ and $\grad\to i\vb p$, so the equation becomes $(-E^2 + \vb p^2 + m^2)\phi = 0$, which holds if and only if $E^2 = \vb p^2+m^2$. For $m = 0$ it is the ordinary wave equation, with $E = \abs{\vb p}$: waves that travel at the speed of light.

\sol{qft}{3} Confinement to $L$ gives momentum $p\sim\hbar/L$ and, once relativistic, energy $\sim\hbar c/L$. This exceeds $2m_ec^2$ for $L\lesssim\hbar/2m_ec\approx200\un{fm}$, about half the Compton wavelength. Squeezing the electron harder creates electron--positron pairs, so the number of particles in the box is no longer fixed: the one-particle wavefunction stops making sense, and you need fields.

\sol{qft}{4} For $r>0$, $\nabla^2f = r^{-1}(rf)''$ for a spherically symmetric $f$; with $f = e^{-mr}/r$, $(rf)'' = m^2e^{-mr}$, so $\nabla^2f = m^2f$ and the equation holds. Near $r = 0$, $f\approx1/r$, and $\nabla^2(1/r) = -4\pi\delta^3(\vb r)$, which supplies the source: $(g/4\pi)(-4\pi\delta^3) = -g\delta^3$. Restoring units turns $m$ in the exponent into $mc/\hbar$, so the range is $\hbar/mc$; for $140\un{MeV}$, $197/140 = 1.4\un{fm}$.

\sol{qft}{5} $\alpha/\pi\approx0.0023$: one-loop corrections are a quarter of a percent, two-loop corrections a few parts per million. Each order is about $400$ times smaller than the last, so a few orders give extraordinary precision. The electron's magnetic moment has been computed through five loops, a correction of order $(\alpha/\pi)^5\sim10^{-13}$, and agrees with experiment to about one part in $10^{12}$.

\sol{qft}{6} $P = \pi^2\hbar c/240d^4$: $1.3\times10^{-3}\un{Pa}$ at $1\un{\mu m}$, and $1.3\times10^{5}\un{Pa}$ --- more than an atmosphere --- at $10\un{nm}$. It equals one atmosphere at $d\approx11\un{nm}$. Micromachines have surfaces separated by such gaps, and the Casimir and van der Waals forces can make moving parts stick together permanently (\emph{stiction}).

\sol{qft}{7} $1/\alpha(91\un{GeV}) = 137.04 - (2/3\pi)\ln(91\,000/0.511) = 137.04 - 2.57 = 134.5$. Heavier charged particles screen too, and bring it to $128$. A charge polarizes the vacuum around it: virtual electron--positron pairs orient themselves, positrons pushed away and electrons pulled toward it (for a negative charge, the reverse), partly screening it, as a dielectric screens a charge. From far away, you see the screened charge, $\alpha\approx1/137$. Probing at shorter distances, inside the cloud, you see more of the bare charge.

\sol{qft}{8} With $\psi' = e^{iq\chi/\hbar}\psi$ and $\vb A' = \vb A+\grad\chi$,
\[
  (-i\hbar\grad - q\vb A')\psi' = e^{iq\chi/\hbar}\big(-i\hbar\grad\psi + q\grad\chi\,\psi - q\vb A\psi - q\grad\chi\,\psi\big) = e^{iq\chi/\hbar}(-i\hbar\grad - q\vb A)\psi.
\]
Applying this twice, the kinetic term transforms into $e^{iq\chi/\hbar}$ times its old value. On the left, $i\hbar\partial_t\psi' = e^{iq\chi/\hbar}(i\hbar\partial_t - q\,\partial_t\chi)\psi$, and the extra term is cancelled by the change in $q\phi\psi$. Every term picks up the same phase, so the equation is unchanged. Without the potentials, the derivatives of $\chi$ would be left over: a phase that varies from place to place would change the physics. Demanding that it not --- \emph{local} phase invariance --- forces you to introduce $\vb A$ and $\phi$, coupled in exactly this way. The electromagnetic field is the price of a local symmetry.

\sol{qft}{9} $V$ depends only on $x = \abs\phi^2$: $\dd V/\dd x = -\mu^2+2\lambda x = 0$ at $\abs\phi = \mu/\sqrt{2\lambda}$, i.e.\ $v = \mu/\sqrt\lambda$. Substituting, $V = -\tfrac12\mu^2(v+h)^2 + \tfrac14\lambda(v+h)^4$, independent of $\theta$. Then $\partial^2V/\partial h^2 = -\mu^2 + 3\lambda v^2 = 2\mu^2$ at $h = 0$: a particle of mass $\sqrt2\mu$ (the Higgs). $\theta$ does not appear in $V$ at all, so moving along the circle of minima costs no energy: a massless Goldstone boson. In a gauge theory, as in the Standard Model, the Goldstone boson becomes the extra polarization of a massive $W$ or $Z$.

\sol{qft}{10} $E_P = 1.2\times10^{19}\un{GeV} = 2.0\times10^9\un{J}$ and $\hbar c = 3.2\times10^{-26}\un{J\,m}$, so $E_P^4/(\hbar c)^3\approx5\times10^{113}\un{J/m^3}$. That is $10^{123}$ times the observed value. A $1\un{TeV}$ cutoff gives $2\times10^{49}\un{J/m^3}$, still $10^{59}$ times too big. Either some unknown principle cancels these zero-point energies to $123$ (or at least $59$) decimal places, leaving a tiny remainder, or we do not understand how the vacuum gravitates. This is the cosmological constant problem.

\sol{qft}{11} In natural units $G_F$ has dimensions of $\text{energy}^{-2}$ and a cross-section has dimensions of $\text{energy}^{-2}$, so $\sigma\sim G_F^2E^2$. This exceeds the unitarity limit $\sim1/E^2$ when $G_F^2E^4\gtrsim1$, at $E\sim G_F^{-1/2}\approx300\un{GeV}$. Something new must appear by then. Solving $G_F = \sqrt2g^2/8m_W^2$ gives $m_W = \sqrt{\sqrt2g^2/8G_F} = 80\un{GeV}$, matching the measured $80.4\un{GeV}$. Fermi's contact interaction is the low-energy limit of $W$ exchange, just as a pointlike Coulomb interaction would be the low-energy limit of a massless photon's.
